\section{Critical zeros constrain every fixed smoothing order (NS-83)}
\label{sec:ns83-critical-budget}

\paragraph{Purpose and attribution.}
Burnol's critical-zero dual vectors strengthen the classical
Nyman--Beurling distance lower bound by counting squared multiplicities
\cite[Sections 3--5]{BurnolLower2002}. We adapt that argument to each
fixed number of Hardy integrations. The logarithmic obstruction survives;
its constant changes. This is an application of the classical mechanism,
not a new zero theorem or the still-missing upper estimate.

For an integer $q\geq1$ fixed independently of $N$, put
\[
 \tau_q=J^{q-1}\chi,\qquad \sigma_{q,n}=J^{q-1}\rho_n,
 \qquad E_N^{(q)}=\inf_{c\in\mathbb R^N}
 \left\|\tau_q-\sum_{n\leq N}c_n\sigma_{q,n}\right\|_2^2.
\]
Thus $E_N^{(2)}=E_N$ is the present smoothed squared error.
Let $\mathcal Z_c$ be the set of distinct zeta zeros on $\Re s=1/2$,
including both signs of the ordinate; write $m_\rho$ for multiplicity.

\begin{proposition}[Fixed-order critical-zero lower bound]
\label{prop:ns83-log-floor}
For every fixed integer $q\geq1$,
\begin{equation}
 \liminf_{N\to\infty}(\log N)E_N^{(q)}
 \geq C_q:=\sum_{\rho\in\mathcal Z_c}
                  \frac{m_\rho^2}{|\rho|^{2q}}>0.
 \label{eq:ns83-log-floor}
\end{equation}
The series is finite. More precisely, every finite subset of
$\mathcal Z_c$ supplies the corresponding finite-sum lower bound.
This is an asymptotic assertion with no numerical onset threshold.
It does not assert $E_N^{(q)}\geq C_q/\log N$ at any specified finite $N$.
\end{proposition}

\subsection{The complete dual-vector construction}
Work first in the complexification of $L^2(0,\infty;dx)$, with inner
product linear in the first argument. For real target and atoms,
complex coefficients cannot improve the infimum: the squared norm is
the sum of the real-residual and imaginary-combination squared norms.
Use the Mellin convention $\widehat f(s)=\int_0^\infty f(x)x^{s-1}dx$.
Let
\[
 (Mf)(x)=x^{-1}\int_0^x f(t)dt,\qquad C=M-I,\qquad W=I-J.
\]
Here $M=J^*$; $C$ and $W$ are unitary, with boundary multipliers
$s/(1-s)$ and $(s-1)/s$ respectively. All commute with the normalized
dilations $D_\theta f(x)=\theta^{-1/2}f(x/\theta)$.
The transformed first atom $a_q=W\sigma_{q,1}$ and target $h_q=W\tau_q$
are supported in $(0,1]$ and have transforms
\begin{equation}
 \widehat a_q(s)=\frac{(1-s)\zeta(s)}{s^{q+1}},\qquad
 T_q(s):=\widehat h_q(s)=\frac{s-1}{s^{q+1}}.
 \label{eq:ns83-transforms}
\end{equation}
Indeed the full exterior $1/x$ tail of each atom is cancelled by $W$,
while $J$ preserves support in $(0,1]$. Explicitly, with $L_x=\log(1/x)$,
\[
 h_q(x)=\left(\frac{L_x^{q-1}}{(q-1)!}
             -\frac{L_x^q}{q!}\right)\chi(x).
\]
Thus no tail-space identification or boundary regularity is presumed.

Write
\[
 \gamma_+(s)=\pi^{1/2-s}\frac{\Gamma(s/2)}{\Gamma((1-s)/2)},
 \quad U(s)=\frac{s}{1-s}\gamma_+(s),
 \quad V_q(s)=\left(\frac{s}{1-s}\right)^{q+2}\gamma_+(s).
\]
These functions define unitary Mellin multipliers on the critical line.
The functional equation identifies $\gamma_+(s)$ with
$\zeta(1-s)/\zeta(s)$ by analytic continuation; the gamma expression is
regular at the critical zeros. We never divide two numerical zero values.
Since $a_q$ is real and
$\widehat a_q(1-s)=V_q(s)\widehat a_q(s)$, one has
\[
 V_q a_q=\mathcal I a_q,\qquad
 (\mathcal I f)(x)=x^{-1}f(1/x).
\]
In particular $V_qD_\theta a_q$ is supported in $[\theta,\infty)$.
Let $Q_\lambda$ be multiplication by $\mathbf1_{[\lambda,\infty)}$.

We spell out the analytic input from Burnol's construction to check that
extra integrations do not change its domain. For
$\psi_{z,k}(x)=\log^k(1/x)x^{-z}\chi(x)$, initially $\Re z<1/2$,
his oscillatory-integral calculation
\cite[Section 4, especially the expansion preceding Theorem 4.11]{BurnolLower2002}
gives an extension $\Phi_{z,k}=U\psi_{z,k}$ to $0<\Re z<1$ with
\[
 \Phi_{z,k}(x)=\partial_z^k\{U(z)x^{-z}\}+O(x^2)
                   \quad(0<x\leq1),\qquad
 \Phi_{z,k}(x)=O(x^{-1})\quad(x\geq1).
\]
For any fixed $k$ these bounds are locally uniform in $z$ in the strip.
The small-$x$ error follows directly from his convergent even-power
series; the large-$x$ bound follows from integration by parts in its
oscillatory integral. No RH assumption enters these estimates.

Here $V_q=C^{q+1}U$. Applying $M$ to $x^{-z}$ multiplies it by
$1/(1-z)$, and applying $M$ to $O(x^2)$ leaves $O(x^2)$ at zero.
At infinity it sends $O((1+\log x)^r/x)$ to
$O((1+\log x)^{r+1}/x)$, with the integral over $(0,1)$ finite.
Differentiation in $z$ is legitimate on compact substrips. Therefore
\begin{equation}
 V_q\psi_{z,k}(x)=\partial_z^k\{V_q(z)x^{-z}\}+O(x^2)
                      \quad(0<x\leq1),\qquad
 V_q\psi_{z,k}(x)=O((1+\log x)^{q+1}/x)\quad(x\geq1).
 \label{eq:ns83-dual-estimates}
\end{equation}
The expressions before truncation are generalized functions outside their
initial $L^2$ domain; only the truncated objects below are asserted to be
Hilbert-space vectors. Both the complete infinite tail and the singular
endpoint have now been controlled.

For $\Re w>1/2$ put
\[
 k_{w,k}(x)=\log^k(1/x)x^{\overline w-1}\chi(x),\qquad
 Y^{\lambda,q}_{w,k}=V_q^{-1}Q_\lambda V_q k_{w,k}.
\]
Equation~\eqref{eq:ns83-dual-estimates}, with $z=1-\overline w$,
proves that $Y^{\lambda,q}_{w,k}$ has an $L^2$ limit as $w$ approaches
any fixed critical-line point $\rho$ from the right. Denote it by
$Y^{\lambda,q}_{\rho,k}$. The estimates are uniform near that point,
so this is norm convergence, not a formal boundary evaluation.

If $\lambda\leq\theta\leq1$, the support identity above gives, in
the initial domain,
\[
 \langle D_\theta a_q,Y^{\lambda,q}_{w,k}\rangle
 =\langle D_\theta a_q,k_{w,k}\rangle
 =(-1)^k\partial_w^k\{\theta^{w-1/2}\widehat a_q(w)\}.
\]
Pass to the norm limit at $w=\rho$. For $k<m_\rho$ the right side
vanishes. Thus every such vector is exactly orthogonal to all
transformed atoms whose dilation parameter lies in $[\lambda,1]$.
In particular it is orthogonal to the first $N$ integer atoms when
$\lambda=1/N$; scalar dilation normalizations do not change the span.

Put $L=\log(1/\lambda)$ and
$X^{\lambda,q}_{\rho,k}=L^{-k-1/2}Y^{\lambda,q}_{\rho,k}$.
For a fixed finite set of zeros, the Gram matrix of these vectors tends
to a block diagonal matrix with blocks
\begin{equation}
 \mathcal H_{m_\rho}=\left(\frac1{k+l+1}\right)_{0\leq k,l<m_\rho}.
 \label{eq:ns83-hilbert-blocks}
\end{equation}
To see this, use unitarity of $V_q$. On $[\lambda,1]$ the leading
term of $V_q k_{\rho,k}$ is
$V_q(\rho)x^{\overline\rho-1}\log^k(1/x)$, of unit phase.
For the same zero its squared-pair integral is
$L^{k+l+1}/(k+l+1)$. Distinct zeros give an oscillatory integral
$O(L^{k+l})$. Derivatives of the multiplier give lower logarithmic
degrees; the $O(x^2)$ error and complete tail contribute bounded terms.
After the displayed normalization all these other terms vanish.

The target pairings require a separate endpoint check. Since
$h_q=J^{q-1}h_1$ and $V_q=C^{q-1}V_1$, while $CJ=M$, one has
\[
 V_qh_q=M^{q-1}V_1h_1
       =M^{q-1}(I-M)\frac{\sin(2\pi x)}{\pi x}.
\]
The last identity is Burnol's target calculation
\cite[Theorem 5.4]{BurnolLower2002}. It is $O(x^2)$ at zero and
$O((1+\log x)^{q-1}/x)$ at infinity. These bounds and
\eqref{eq:ns83-dual-estimates} give an integrable majorant for the
pairings, uniform as $w\to\rho$ and $\lambda\downarrow0$.
Consequently
\[
 \lim_{\lambda\downarrow0}\langle h_q,Y^{\lambda,q}_{\rho,k}\rangle
       =(-1)^k T_q^{(k)}(\rho),
\]
and hence
\begin{equation}
 \sqrt L\,\langle h_q,X^{\lambda,q}_{\rho,k}\rangle
 \longrightarrow
 \begin{cases}T_q(\rho),&k=0,\\0,&k\geq1.\end{cases}
 \label{eq:ns83-target-pairings}
\end{equation}

\begin{proof}[Proof of Proposition~\ref{prop:ns83-log-floor}]
Fix a finite subset of critical zeros and include $0\leq k<m_\rho$
for each. The squared error is bounded below by the squared norm of the
projection of $h_q$ onto their dual-vector span. The limiting Gram in
\eqref{eq:ns83-hilbert-blocks} is positive definite, so its inverse
converges. The elementary Cauchy-matrix formula gives
$(\mathcal H_m^{-1})_{00}=m^2$. Using
\eqref{eq:ns83-target-pairings}, the limiting lower bound after
multiplication by $L$ is
\[
 \sum_\rho m_\rho^2|T_q(\rho)|^2
 =\sum_\rho\frac{m_\rho^2}{|\rho|^{2q}},
\]
because $|\rho-1|=|\rho|$ on the critical line.
Take $\lambda=1/N$ and then the supremum over finite zero sets.
The classical zero-count estimate (for example \cite{HasanalizadeShenWong2021}) gives $O(\log(2+T))$ zeros with
multiplicity in each unit-height interval and therefore the same bound
on an individual multiplicity. Comparison with
$\int_1^\infty \log^2(2+t)t^{-2q}dt$ proves convergence of the series.
The existence of a critical zero makes it strictly positive. The entire
argument uses actual critical zeros and no hypothesis about other zeros.
\end{proof}

\subsection{What relative contraction can still be sought}
For the present $q=2$ problem set
\[
 \gamma_j=1-\frac{E_{2^{j+1}}}{E_{2^j}},\qquad j\geq1.
\]
Finite errors are positive (an exact finite identity would contradict a
critical zero of $\zeta$), and the nested projection spaces give
$0\leq\gamma_j<1$. For any fixed starting index $j_0$, telescoping and
Proposition~\ref{prop:ns83-log-floor} imply
\begin{equation}
 \sum_{j=j_0}^{K-1}-\log(1-\gamma_j)
 =\log\frac{E_{2^{j_0}}}{E_{2^K}}
 \leq\log K+O(1),\qquad
 \sum_{j=j_0}^{K-1}\gamma_j\leq\log K+O(1).
 \label{eq:ns83-gain-budget}
\end{equation}
In particular no fixed positive fraction of the entire error can be
removed at every sufficiently large doubling. Nor can one have
$\gamma_j\geq a/j$ eventually for a fixed $a>1$.
The same lower bound also shows that $E_N^{(q)}=O((\log N)^{-1-\varepsilon})$ is impossible
for any fixed $q$ and $\varepsilon>0$; a power of $N$ is still faster.
These statements do not give a pointwise bound $\gamma_j\leq1/j$:
individual large gains and irregular blocks are allowed.

A sufficient arithmetic estimate of the scale $a/j$ with $0<a\leq1$,
or a weaker nonsummable scale, is not excluded. None is proved here.
NS-80/82's uniform capture of \emph{available block gain} is fully
consistent with this restriction on gain relative to the entire error.
The calculation does not prove $E_N\to0$, an asymptotic equality, the
original Weil floor, G2 or RH. The fixed-order qualification is essential;
letting the smoothing order change with $N$ changes the target and norm.

\paragraph{Certified constant, not a finite error bound.}
At 256 and 384 bits, the first ten positive critical zeros and their
conjugates give the partial-sum lower bound
\[
 C_2\geq2\sum_{j=1}^{10}(t_j^2+1/4)^{-2}
 >7.2332433\cdot10^{-5}.
\]
Here $t_j$ denotes the $j$th positive zero ordinate. Each zero is
counted with weight one, so unknown
larger multiplicities could only increase the constant. All omitted
critical-zero terms are nonnegative: this is a lower bound from a finite
subset, not an approximation to the full sum with a missing tail.
There is no numerical bound on the onset of
\eqref{eq:ns83-log-floor}; the constant is not used as a certified bound
on the current size-256 error.
